\section{Transport Invariant and Signed Lift}

We now identify the binary transport invariant governing the lift
\[
\Gthirty \to \Gfifteen.
\]
The key point is that the sign on a core edge is not assigned ad hoc:
it is induced by the lifted transport and admits a canonical
geometric interpretation on the dodecahedron. The resulting structure
fits naturally into the theory of signed graphs and $2$-lifts
(see \cite{zaslavsky}).

\subsection{The antipodal lift}

The graph $\Gthirty$ is a $2$-fold cover of $\Gfifteen$. Thus each vertex
$x \in V(\Gfifteen)$ has a two-point fiber
\[
\pi^{-1}(x)=\{x^+,x^-\}\subset V(\Gthirty).
\]
For an edge $e=\{x,y\}\in E(\Gfifteen)$, the adjacency upstairs connects the
two fibers $\pi^{-1}(x)$ and $\pi^{-1}(y)$ in exactly one of two ways:
either sheet is preserved, or sheet is exchanged. This binary distinction
defines the transport parity of the edge.

\subsection{Canonical dodecahedral model of the core}

By the identification
\[
\Gfifteen \cong L(\mathrm{Petersen}),
\]
we may regard the vertices of $\Gfifteen$ as the $15$ opposite-edge classes
of the dodecahedron. This identification is classical (see
\cite{godsil_royle}).

Each vertex of $\Gthirty$ then corresponds to one of
the two actual dodecahedral edges in the corresponding opposite-edge class.

Let $e,f$ be two dodecahedral edges arising from adjacent vertices of
$\Gfifteen$. Their relation is classified as follows:

\begin{itemize}
\item \emph{local} if $e$ and $f$ share a vertex and lie on a common face;
\item \emph{nonlocal} if they neither share a vertex nor lie on a common face,
and are not opposite edges.
\end{itemize}

\begin{proposition}[Uniformity of lifted edge type]
Let $e=\{x,y\}\in E(\Gfifteen)$. Under the identification with opposite-edge
classes of the dodecahedron, the two lifted edges in $\Gthirty$ above $e$
are either both local or both nonlocal. In particular, the only possible
patterns are $(LL)$ and $(NN)$.
\end{proposition}

\begin{proof}
Each vertex of $\Gfifteen$ corresponds to a pair of opposite edges of the
dodecahedron, and the fiber over that vertex consists of the two actual
edges in the pair.

Adjacency in $\Gfifteen$ is determined by a fixed combinatorial relation
between opposite-edge classes. Lifting this relation to $\Gthirty$ replaces
each class by its two representatives but does not alter the underlying
incidence geometry.

Since the local/nonlocal classification depends only on the geometric
relation between the underlying edges (sharing a vertex and face, or not),
and this relation is preserved under passage to representatives within
opposite-edge classes, the two lifted edges must have the same type.

Thus only $(LL)$ and $(NN)$ can occur.
\end{proof}

\subsection{Definition of the Thalean cocycle}

\begin{definition}[Thalean cocycle]
Let $e=\{x,y\}\in E(\Gfifteen)$. Consider the two lifted edges in
$\Gthirty$ connecting the fibers $\pi^{-1}(x)$ and $\pi^{-1}(y)$, and
interpret them via the identification of $\Gfifteen$ with the graph of
opposite-edge classes of the dodecahedron. Define
\[
\varepsilon(e)\in \ZZtwo
\]
by
\[
\varepsilon(e)=
\begin{cases}
0,&\text{if the lifted pattern is }(LL),\\
1,&\text{if the lifted pattern is }(NN).
\end{cases}
\]
Equivalently, $\varepsilon(e)=0$ records trivial sheet holonomy across
the edge, while $\varepsilon(e)=1$ records nontrivial sheet exchange.
This $\ZZtwo$-valued edge function is called the \emph{Thalean cocycle}.
\end{definition}

\subsection{Cocycle theorem}

\begin{theorem}[Thalean cocycle theorem]
The edge function
\[
\varepsilon : E(\Gfifteen) \to \ZZtwo
\]
defined above has the following properties:

\begin{enumerate}[label=(\roman*)]
\item It is canonically induced by the lifted transport
$\Gthirty\to \Gfifteen$.

\item It records whether transport across a core edge preserves or swaps
sheet in the double cover.

\item Its cycle parity computes the sheet holonomy of lifted loops:
for any cycle $C\subset \Gfifteen$,
\[
\omega(C):=\sum_{e\in C}\varepsilon(e)\pmod 2
\]
is trivial if and only if the lifted loop closes on the same sheet.

\item It defines a nontrivial cohomology class in
$H^1(\Gfifteen,\ZZtwo)$.
\end{enumerate}
\end{theorem}

\begin{proof}
The lift $\Gthirty\to \Gfifteen$ provides a canonical two-sheet structure
over each vertex of $\Gfifteen$. For each base edge, the induced adjacency
between the two fibers is either sheet-preserving or sheet-swapping. This
binary distinction is invariant under exchanging the two labels in a fiber,
so it defines a well-defined element of $\ZZtwo$ on each edge.

Under the canonical identification of $\Gfifteen$ with the graph of
opposite edge-pairs of the dodecahedron, the two lifted realizations of
each base edge are simultaneously local or simultaneously nonlocal,
yielding exactly the patterns $(LL)$ and $(NN)$.

For a cycle $C$ in $\Gfifteen$, the parity sum of edge values records
whether the lifted walk in $\Gthirty$ returns to its original sheet or
to the opposite sheet. In particular, the cocycle is nontrivial if and
only if there exists a cycle with $\omega(C)=1$.

By the transport construction, there exist cycles with both trivial and
nontrivial holonomy. Hence the cocycle represents a nontrivial class in
$H^1(\Gfifteen,\ZZtwo)$.
\end{proof}

\subsection{Holonomy classes of pentagons}

The cocycle does not create the pentagons of $\Gfifteen$; rather, it
polarizes them into trivial and nontrivial holonomy classes. In
particular, the parity sum
\[
\omega(C)=\sum_{e\in C}\varepsilon(e)\pmod 2
\]
distinguishes pentagons whose lifted transport closes trivially from
those carrying nontrivial sheet-flip holonomy. In particular, the
existence of both trivial and nontrivial holonomy classes provides a
direct witness of the nontriviality of the cocycle.

\subsection{Signed-lift interpretation}

If one passes from $\ZZtwo=\{0,1\}$ to $\{\pm1\}$ by the rule
\[
\sigma(e)=(-1)^{\varepsilon(e)},
\]
then $\sigma(e)=+1$ corresponds to trivial holonomy and
$\sigma(e)=-1$ to nontrivial holonomy. Thus the Thalean cocycle may be
viewed equivalently as a signed edge function on $\Gfifteen$ or as the
$\ZZtwo$-voltage defining the double cover $\Gthirty\to \Gfifteen$, in
the sense of signed graph theory \cite{zaslavsky}.

This cocycle is the fundamental transport invariant of the Petersen
core. While it does not appear in quadratic sector observables, it is
detected by loop holonomy and governs the topology of lifted transport.
